\section{Conclusion}
\label{sec:conclusion}

This systematic review analyzed 13 studies on wearable seizure detection and forecasting using deep learning approaches. The evidence base comprises 912 participants across nine detection studies and four forecasting studies published between 2013 and 2026. Three key findings emerge from this analysis.

First, detection approaches are approaching clinical readiness. Two studies meet the ILAE Phase 3 benchmark for false positive rates below 0.05 to 0.1 per hour. \textcite{spahrDeepLearningbasedDetection2025} achieved 0.0054/h FPR using an ensemble of 30 CNN models with single-modality ACC data. \textcite{fineDetectionKeyAutomated2025} achieved 0.023/h FPR using a 6-axis ACC and gyroscope band. Accelerometer-based detection shows strong translation potential for convulsive seizures. However, a significant gap remains between EMU validation and home deployment. Only three detection studies include home validation, and FPR increases substantially in real-world settings.

Second, forecasting approaches face a consistent performance ceiling. AUC values plateau around 0.74 to 0.75 across studies using different methods, modalities, and patient populations. \textcite{meiselMachineLearningWristband2020} achieved AUC 0.74 with LOSO validation. \textcite{nasseriAmbulatorySeizureForecasting2021} achieved AUC 0.80 with patient-specific models. \textcite{stirlingForecastingSeizureLikelihood2021} achieved AUC 0.74 with heart rate cycles. This consistency suggests fundamental limits on pre-ictal predictability using non-invasive wearable sensors. Unlike detection, forecasting lacks established clinical benchmarks and regulatory pathways.

Third, three barriers impede clinical translation. The personalization dilemma remains unresolved. Global models enable immediate deployment but achieve only 62\% patient success in LOSO validation. Patient-specific models achieve higher success rates but require weeks to months of individual data, creating a cold start problem. Seizure type coverage is limited to convulsive seizures with motor manifestations. Approximately 30\% of people with epilepsy have non-convulsive seizures that remain undetectable with current wearable approaches. Real-world deployment presents challenges not captured in hospital validation, as evidenced by higher FPR in home settings.

No study in this review has achieved all requirements for widespread clinical adoption. Commercial devices exist for detection including the Embrace watch and NightWatch armband, but no device has regulatory approval for forecasting. Future research should prioritize home validation at scale, standardized benchmarks for forecasting performance, and novel modalities for non-convulsive seizure detection. Mixed personalization approaches that combine population patterns with individual adaptation may offer a path forward.
